\chapter{Generate an ORCA input (menu 3)}
\label{ch:menu3}
\menulabel{3}

\section{Purpose}

Answer a short system profile and get a directly runnable ORCA input file --
method chain, basis set and ECP, starting active space, convergence settings --
plus plain guidance on how to run and what to check afterwards.

\section{What you need}

\begin{itemize}
  \item the system profile: element list, charge, spin multiplicity,
    f-block metal valence, targets
    (energy / geometry / excited / magnetic / spectra);
  \item a structure file (XYZ): generating an input needs a geometry;
  \item optionally the active space (as \code{nel,norb,mult,nroots}; Enter
    skips the \code{\%casscf} block), the convergence tier, and the
    \code{MaxCore} per process in MB (Enter keeps the default 2000; a
    measured $f$-block TRAH-CASSCF needed 9345 MB, so raise it for
    large-basis hard cases -- the guidance quotes the value you chose).
\end{itemize}

\section{Walkthrough}

\lstinputlisting[style=fbkcode, firstline=17, lastline=55,
  title={The menu-3 example session (Ce\textsuperscript{3+}, 4f\textsuperscript{1})}]
  {examples/expected/m3_generate.txt}

The session shows the three ingredients the recipe layer contributes: the
method chain (``CASSCF + SC-NEVPT2''), the basis tiers with their rationales
(the SARC2 tier first for lanthanide wavefunction work; the def2 tier with its
matched 28-core ECP for screening), and the G2 starting active spaces with
their rationale (the 4f-only minimal window, confirmed by a literature
precedent; the 4f+d double shell, provisional). It also shows the convergence
plan and the run guidance.

\section{What you get}

\file{structure.fbk.inp} -- pure ASCII (an assertion, not a hope: ORCA input
parsers break on non-ASCII bytes), directly runnable -- plus the run guidance
printed on screen: method chain, basis, convergence discipline, refusals, and
what to check afterwards.

\section{How to read the numbers}

\begin{readbox}
\begin{itemize}
  \item Tiers are ordered, not ranked: ``1'' is the recommended default for
    the stated profile; the notes say \emph{why} each tier exists (screening,
    magnetic-property chains, correlation energies).
  \item Active-space suggestions are \emph{suggestions}: the electron count
    comes from a standard filling formula, and the orbital selection still
    needs your chemical insight. ``confirmed'' means ``a literature precedent
    exists'', not ``the program knows your ground state''.
  \item Every recommendation carries its provenance; the run guidance repeats
    the refusals that matter (for example: TRAH without its /C auxiliary basis
    is refused rather than silently written).
\end{itemize}
\end{readbox}

\subsection{The relativistic tier line}

When the profile names an $f$-block element or a spin-orbit target
(\code{magnetic}, \code{spectra}), the run guidance adds the relativistic tier
the need implies -- the last line of the listing above.  The tiers come from the
X2C hierarchy for two-component work, and the rule's point is \emph{avoiding
over-configuration}: scalar work uses the one-electron variants, light
main-group spin-orbit work wants the two-electron correction (it can carry a
quarter of the splitting in O$_2$), and the f block is the case where paying for
that correction buys only $1.6$--$4.9$~cm$^{-1}$ on Nd\textsuperscript{3+}
splittings -- so the atomic-mean-field tier is the default there and a higher
tier must be justified by the target accuracy.  One option is ruled out in
every row: the atomic-mean-field variant \code{amfX2C} has a numerically
unstable single-centre scalar approximation.

For f-block spin-orbit work the same source supplies the two-component
SO-CASSCF template this toolkit records as guidance (ORCA itself offers the
one-electron variants only): the whole $4f$ shell, the ground-term manifold
state-averaged in full -- the window is computed as $(2S+1)(2L+1)$ from the
Hund term, giving the source's 52 states for Nd\textsuperscript{3+} -- spin-orbit
contracted bases for the heavy elements, Cholesky threshold $10^{-7}$~Hartree,
and the report of Kramers doublet levels plus each multiplet's centre of
gravity.

\section{Worked example}

The generated input of this session was run for real on ORCA 6.1.1 and its
output is one of the fixture files (\file{generated\_ce3\_sarc2.out}); feeding
that output back into \menu{1} closes the loop. That closed loop is also a
regression fixture of the project: the input the recipe layer writes today must
run and parse.

\section{Caveats, refusals and related sections}

\begin{refusalbox}
\begin{itemize}
  \item The program writes the input file and the run guidance; the engine step
    itself stays with the user.
  \item Unsuitable requests are refused with the reason, not softened: for
    example, a 5f-in-core pseudopotential for f--f transitions is listed as
    not applicable (Chapter~\ref{ch:menu4}), and a geometry-optimisation
    target never gets a correlated layer without an analytic gradient.
  \item The generated \code{\%casscf} block is written only when it is
    complete (\code{nel,norb,mult,nroots}); an incomplete request downgrades to
    a file without that block rather than to a file that will not run.
\end{itemize}
\end{refusalbox}

Related: \menu{4} (the same basis knowledge without generating a file),
Chapter~\ref{ch:criteria} for the convergence rules, and the OpenMolcas
magnetic-property template (generated by the recipe layer; see
Chapter~\ref{ch:development} for where it lives).
